\documentclass[11pt,a4paper]{article}
\usepackage[margin=2.3cm]{geometry}
\usepackage{amsmath,amssymb,amsthm}
\usepackage{booktabs}
\usepackage{array}
\usepackage{longtable}
\usepackage{enumitem}
\usepackage{microtype}
\usepackage{caption}
\usepackage{float}
\usepackage[hidelinks]{hyperref}
\usepackage{xcolor}
\usepackage{tcolorbox}

\captionsetup{font=small,labelfont=bf}
\setlength{\parindent}{0pt}
\setlength{\parskip}{5pt}

\newtheorem{assumption}{Assumption}
\newcommand{\Tr}{\mathcal{T}}
\newcommand{\Dn}{\mathcal{D}}
\newcommand{\ind}{\mathbf{1}}
\newcommand{\lp}{log-\%}
\newcommand{\E}{\mathbb{E}}
\newtcolorbox{keybox}[1]{colback=blue!4!white,colframe=blue!45!black,title={#1},fonttitle=\bfseries}
\newtcolorbox{plainbox}{colback=gray!5!white,colframe=gray!55!black,left=5pt,right=5pt,top=3pt,bottom=3pt}

\title{\textbf{Hurricane Helene and County Labor Markets in Western North Carolina}\\[4pt]
\large Complete Methodological Annotation}
\author{Piergiulio Fasciani \\ \small Bocconi University}
\date{18 November 2026 \\ \small Specification of 18 November 2026; repository \texttt{Helene\_Definitive\_Joint\_Research}}

\begin{document}
\maketitle

\begin{keybox}{What this design does, in three sentences}
Thirty-nine North Carolina counties with verified hurricane damage are compared with a weighted
average of 1{,}150 untreated counties in distant states, where the weights are chosen so that the
comparison tracked the treated counties' pre-storm \emph{movements} in four labor outcomes over
thirty-two months. The treated group is then split into terciles of storm damage, so the paper
estimates not one average effect but a dose-response: the least-damaged third shows no effect at any
horizon, while the most-damaged third loses 4.3\% of resident employment in the first quarter.
Migration enters as an outcome of the same design, never as a control, and shows that the most-damaged
counties' population growth falls well below what the comparison implies in the year after the storm,
which is why their labor force does not recover when their employment does.
\end{keybox}

\tableofcontents

\section{Question, estimand, and what the design can deliver}

\subsection{The two questions}
\begin{enumerate}[nosep,leftmargin=1.4em]
\item \textbf{Does damage intensity drive labor-market losses and their persistence?} Answered
quasi-experimentally, by estimating a separate effect for each damage tercile against never-treated
counties.
\item \textbf{What drives the recovery?} Answered by decomposition --- employment, labor force,
unemployment and migration are each estimated as outcomes of the same design, so the channels can be
read off against one another rather than inferred.
\end{enumerate}

\subsection{Estimand}
Let $Y_{it}$ be an outcome for county $i$ in month $t$, $\Tr$ the set of 39 treated counties, and
$\bar Y_{\Tr t}=\tfrac{1}{39}\sum_{i\in\Tr}Y_{it}$ the equal-weighted treated mean. For a post-event
window $W$,
\begin{equation}
\tau_W=\frac{1}{|W|}\sum_{t\in W}\Big(\bar Y_{\Tr t}(1)-\bar Y_{\Tr t}(0)\Big),
\end{equation}
where $\bar Y_{\Tr t}(0)$ is the untreated potential outcome. For the dose design the same object is
computed within each exposure group $g$,
\begin{equation}
\tau_W(g)=\frac{1}{|W|}\sum_{t\in W}\Big(\bar Y_{g t}(1)-\bar Y_{g t}(0)\Big),\qquad g\in\{\text{low},\text{mid},\text{high}\},
\end{equation}
each with 13 counties. Three design choices are embedded and should be stated rather than discovered:

\begin{itemize}[nosep,leftmargin=1.4em]
\item \textbf{Equal county weights, not population weights.} The estimand is ``the typical affected
county'', not ``the affected region in aggregate''. Buncombe (Asheville, $\approx$261k residents)
counts the same as Alleghany ($\approx$11k). A population-weighted estimand would be close to a study
of Asheville alone.
\item \textbf{Effects on the treated aggregate}, not county-level effects. Weights are fitted to the
group mean, which is far less noisy than any individual rural county.
\item \textbf{Observational treatment.} ``Treated'' means ``recorded positive verified FEMA loss'',
an ex-post exposure definition. Nothing here is randomised.
\end{itemize}

\subsection{Windows}
Three post-event windows are fixed in advance of estimation:
\textbf{immediate} (October--December 2024, 3 months), \textbf{2025} (January--September 2025, 9
months), and \textbf{late} (November 2025--March 2026, 5 months). September 2024 is the storm month
and is excluded as a transition period; October 2025 is excluded because LAUS is missing at source.

\section{Data and outcome construction}

\subsection{Sources}
\begin{center}\small
\begin{tabular}{>{\raggedright}p{3.1cm}>{\raggedright}p{7.2cm}>{\raggedright\arraybackslash}p{4.9cm}}
\toprule
Source & Content and construction & Role \\
\midrule
BLS LAUS (county, monthly) & Resident employment, unemployment, labor force, unemployment rate. Obtained through a documented third-party distributor of the BLS series; the official flat-file route is implemented but returns HTTP 403. & Three of four fitting outcomes; residence-based.\\
BLS QCEW (county, quarterly files carrying monthly employment) & Workplace jobs, all industries, 2022Q1--2026Q1. & Fourth fitting outcome; workplace-based.\\
NOAA nClimGrid-Daily & County precipitation summed over 25--30 September 2024; all six daily values required. & Exposure measure 2 (rainfall).\\
FEMA IHP valid registrations, disasters 4827--4832 & Verified real- and personal-property losses, aggregated to county. & Treatment definition and exposure measure 1.\\
Census PEP 2024; ACS 2023 5-year & Population; households and median home value. & Denominators for the damage index.\\
Census Vintage 2025 components & Annual July-to-July population change and net domestic migration. & Migration outcome.\\
\bottomrule
\end{tabular}
\end{center}

The constructed panel is balanced on the calendar: 3{,}221 counties $\times$ 57 months (January
2022--September 2026) $=183{,}597$ rows. Balance on the calendar is not balance on observation:
QCEW ends in March 2026 and October 2025 LAUS is missing at source. Missing values are kept missing
and never imputed or set to zero.

\subsection{The four fitting outcomes}
\begin{equation}
\begin{aligned}
y^{(1)}_{it}&=100\ln(\text{resident employment}), &\quad y^{(2)}_{it}&=100\ln(\text{resident labor force}),\\
y^{(3)}_{it}&=100\ln(\text{workplace jobs}), &\quad y^{(4)}_{it}&=\text{unemployment rate (pp)}.
\end{aligned}
\end{equation}
The factor $100\ln$ makes a unit change approximately one percent, written \lp\ throughout. Each
county series is first centred on its own 2022 mean,
$\tilde y^{(k)}_{it}=y^{(k)}_{it}-\tfrac{1}{12}\sum_{s\in 2022}y^{(k)}_{is}$, so that counties of very
different size are comparable in movements.

Two derived outcomes are reported alongside: $y^{(3)}-y^{(1)}$, the log ratio of workplace jobs to
resident employment (where the jobs are versus where the workers live), and $y^{(1)}-y^{(2)}$, the log
employment-to-labor-force ratio (approximately the negative of the unemployment rate). Both are exact
linear combinations of the first three, which matters for inference and is discussed in
Section~\ref{sec:multiplicity}.

\begin{plainbox}
\textbf{Why four outcomes rather than one.} A synthetic control fitted to a single series can match it
by coincidence while mismatching everything else. Requiring one weight vector to track employment,
labor force, workplace jobs and the unemployment rate simultaneously over-identifies the weights.
It also makes the two ratio outcomes internally consistent, because the same control is subtracted
from numerator and denominator.
\end{plainbox}

\section{Treatment group}

\subsection{Definition}
A North Carolina county is treated if it appears in the NOAA eleven-state rainfall sample and records
strictly positive verified FEMA losses (real plus personal property) under disasters 4827--4832. This
yields \textbf{39 counties}, essentially the western third of the state.

The roster is \emph{frozen} in two independent reference files. The source build regenerates the
FEMA-positive definition from raw records at every run and halts if the regenerated set differs from
the frozen roster. This protects against two distinct hazards: later FEMA revisions silently changing
the sample, and post-hoc reclassification after outcomes have been seen.

\subsection{Exposure within the treated group}
Two continuous measures are carried for the dose design.

\textbf{Rainfall.} Six-day total, 25--30 September 2024, millimetres. Across the 39: mean 214,
SD 95, range 70--402.

\textbf{HDSI} (Helene damage severity index), a composite:
\begin{equation}
\text{HDSI}=0.4\,R_{\text{rain}}+0.6\,R_{\text{burden}},
\end{equation}
where $R$ are percentile-rank scores on $[0,100]$ within the complete-case sample and economic burden
is $100\times(\text{verified loss})/(\text{households}\times\text{median home value})$. Across the 39:
mean 62.8, SD 22.4, range 25.1--98.1.

HDSI is descriptive. It does not define treatment and is not claimed exogenous: damage correlates with
terrain, income and industry mix. Rainfall is included precisely because it is pure meteorology with no
FEMA claim behaviour in it, and therefore provides a check on whether results are an artefact of how
the composite is built.

\subsection{What ``treated'' does not mean}
Treated counties are Appalachian, lower-income, older and more tourism-dependent than the state
average. The design handles this by matching on pre-event \emph{dynamics} (Section~\ref{sec:weights}),
not by conditioning on covariates.

\section{Control group}

\subsection{Geographic exclusions}
A county is eligible as a donor only if its state lies outside three exclusion sets, applied by state
FIPS:
\begin{itemize}[nosep,leftmargin=1.4em]
\item \textbf{Impact states} (storm track, flooding, or declared disasters): AL, FL, GA, IL, IN, KY, NC, OH, SC, TN, VA, WV.
\item \textbf{Buffer states} (adjacent to impact states; commuting, migration and supply-chain spillovers): AR, DC, IA, MD, MI, MS, MO, PA, WI.
\item \textbf{Other 2024 hurricane states}: LA, TX (Beryl, Francine).
\end{itemize}
This leaves 1{,}161 eligible counties across 27 states and Puerto Rico. The rationale is SUTVA: a valid
control must be unaffected by the storm, including indirectly through displaced residents,
reconstruction labor demand, or regional supply chains. The cost is comparability --- the donors are
not Appalachian --- and the weighting step is what recovers it.

\subsection{Completeness screen}
Of the 1{,}161, a county is usable only if all four fitting outcomes are observed and strictly positive
in each of the 24 months of 2022--2023, leaving \textbf{1{,}150 donors}. Only those 24 months enter the
screen, so no 2024 or post-storm observation influences who is eligible.

\subsection{The fitted control}
The synthetic control is $\widehat{\bar Y}_{\Tr t}(0)=\sum_{j\in\Dn}\widehat w_j\,\tilde y_{jt}$, with
$\widehat w\ge0$ and $\sum_j\widehat w_j=1$. The fitted vector places positive weight on
\textbf{115 donors}, with effective count $\big(\sum_j\widehat w_j^2\big)^{-1}=63.9$; the largest single
weight is 3.76\% and the ten largest sum to 27.7\%. Weight is distributed across Wyoming, Arizona,
South Dakota, Idaho, Nevada, Utah, Puerto Rico and New Mexico. No donor dominates, which matters for
robustness to any single county's idiosyncrasies.

\section{The weighting problem}\label{sec:weights}

\subsection{Objective}
Let $T_{\text{pre}}$ be the fitting window, $K=4$ outcomes, $J=1{,}150$ donors, $s_k$ a fixed scale per
outcome, $u_j=1/J$ the uniform vector, and $\Delta^J$ the probability simplex. The weights solve
\begin{equation}\label{eq:obj}
\min_{w\in\Delta^{J},\,a\in\mathbb{R}^K}\;
\frac{1}{2KT_{\text{pre}}}\sum_{k=1}^{K}\sum_{t\in T_{\text{pre}}}
\left(\frac{\bar{\tilde y}^{(k)}_{\Tr t}-a_k-\sum_j\tilde y^{(k)}_{jt}w_j}{s_k}\right)^{2}
+\frac{\lambda}{2}\sum_{j=1}^{J}\big(w_j-u_j\big)^2 .
\end{equation}

Three ingredients, each doing specific work:

\begin{description}[leftmargin=1.2em,style=nextline]
\item[The intercepts $a_k$.] Each outcome block is demeaned over the fitting window, which profiles out
$a_k$. The weights therefore match the treated aggregate's \emph{movements}, while a constant
treated-control level difference is permitted and later absorbed by county fixed effects. This is the
synthetic difference-in-differences idea of Arkhangelsky, Athey, Hirshberg, Imbens and Wager (2021):
the control need not sit at the same level, only move in parallel. Fitted intercepts are
$+0.013$, $-0.002$, $-0.042$ and $-0.017$ for the four outcomes.
\item[The scales $s_k$.] $s_k$ is the standard deviation of the treated aggregate's 2022 values of
outcome $k$, floored at 0.25~\lp\ for the three log outcomes and 0.15~pp for the rate. Realised values
are 0.915, 0.804, 1.644 and 0.279. Scaling puts outcomes on a common footing so that none dominates
the fit through its units.
\item[The ridge toward uniform.] The penalty $\tfrac{\lambda}{2}\|w-u\|^2$ shrinks toward equal weights.
With 1{,}150 donors and $4\times32=128$ fitting equations the unpenalised problem is badly
under-determined; the penalty selects the most diffuse weight vector among those fitting well, which
limits over-fitting to pre-period noise. The simplex constraint rules out extrapolation.
\end{description}

\subsection{Choice of $\lambda$}
$\lambda$ is chosen out of sample within the pre-treatment window. Partition $T_{\text{pre}}$ into an
estimation fold (January--December 2022) and a validation fold (January--December 2023). For each
$\lambda\in\Lambda=\{0.03,0.1,0.3,1,3,10\}$, solve \eqref{eq:obj} on the estimation fold alone and
score the same standardised loss on the held-out fold, giving $R(\lambda)$. The selected penalty is
\begin{equation}
\widehat\lambda=\max\Big\{\lambda\in\Lambda:\;R(\lambda)\le(1+\eta)\min_{\lambda'\in\Lambda}R(\lambda')\Big\},\qquad\eta=0.01 .
\end{equation}
Taking the largest rather than the minimising $\lambda$ inside the tolerance band is deliberate: among
weight vectors whose out-of-sample accuracy is practically indistinguishable, the more diffuse one is
preferred, since it spreads mass over more donors and is less exposed to any single comparison county.
This is the familiar one-standard-error heuristic with a fixed relative tolerance.

Realised scores are $0.1543$, $0.1433$, $0.1236$, $0.1094$, $0.1126$, $0.1198$. The minimum is at
$\lambda=1$; the next stronger candidate $\lambda=3$ sits 3.0\% above it, outside the band. The rule
therefore returns the minimiser, $\widehat\lambda=1$: the tolerance is a safeguard that did not bind.

\subsection{Final fit and the backdated fit}
Given $\widehat\lambda=1$, the \textbf{definitive weights} solve \eqref{eq:obj} over the full pre-storm
window, January 2022--August 2024 (32 months). Nothing after August 2024 enters any fit.

A \textbf{backdated} vector solves the identical problem over January 2022--December 2023 only, with
the same $\lambda$ and the same intercept structure. It is never used for estimation. Its sole purpose
is an out-of-sample test: predict January--August 2024 using weights that never saw those months
(Section~\ref{sec:validation}). It places weight on 125 donors, effective count 69.8.

\subsection{Solver}
The problem is a strongly convex quadratic over the simplex, solved in Mata by accelerated projected
gradient (FISTA) with exact Euclidean projection onto the simplex, a Lipschitz step from the leading
eigenvalue, and a duality-gap stopping rule $g^\top w-\min_j g_j\le2\times10^{-5}$ where $g$ is the
gradient. The gap is a certified bound on sub-optimality over the simplex. Convergence is asserted,
not assumed: a run exceeding tolerance halts the pipeline. Realised gap $1.9\times10^{-5}$ in 7{,}000
iterations.

\begin{plainbox}
\textbf{Is this canonical synthetic DiD?} Not entirely, and the paper should say so. Canonical SDID has
three components: unit weights with an intercept and ridge penalty; \emph{time} weights that
down-weight pre-periods unlike the post-period; and a weighted two-way regression. This implementation
has the first and third with \emph{uniform} time weights. A formal SDID time-weight vector was computed
outside the pipeline as a check: it concentrated mass on 2023--2024 and moved the window effects by at
most $\approx0.25$~\lp\ (workplace jobs) and $\le0.05$ for the others. Uniform time weights are retained
for transparency --- every pre-event month counts equally in the reference.
\end{plainbox}

\section{Main event study}\label{sec:es}

\subsection{Specification}
The estimation sample is the 39 treated counties plus the 115 active donors over January 2022--March
2026, excluding September 2024 and October 2025: $154\times49=7{,}546$ county-months. For each outcome
$y$ in raw units (the fixed effects perform the centring),
\begin{equation}\label{eq:twfe}
y_{it}=\alpha_i+\gamma_t+\sum_{k\in\mathcal{K}}\beta_k\,D^{k}_{it}+\varepsilon_{it},
\qquad D^{k}_{it}=\ind\{i\in\Tr\}\cdot\ind\{k_t=k\},
\end{equation}
estimated by weighted least squares with analytic weights $\omega_i=1/39$ for treated counties and
$\omega_i=\widehat w_i$ for donors, county fixed effects absorbed, calendar-month effects included, and
standard errors clustered by county. Event time is $k_t=t-\text{September 2024}$, with indicators for
$k\in\{-32,\dots,18\}\setminus\{-1,0,13\}$: 48 regressors, $k=-1$ omitted to fix the parameterisation,
$k=0$ and $k=13$ absent from the sample. Inference uses $\mathrm{df}=G-1=153$, the cluster count less
one.

\subsection{Reference: the pre-event mean}
Every reported coefficient and window effect is expressed relative to the mean of the 32 pre-storm
months:
\begin{equation}
\tilde\beta_k=\beta_k-\frac{1}{32}\sum_{m=-32}^{-1}\beta_m,\qquad
\widehat\tau_W=\frac{1}{|W|}\sum_{k\in W}\beta_k-\frac{1}{32}\sum_{m=-32}^{-1}\beta_m,
\qquad \beta_{-1}\equiv0,
\end{equation}
with standard errors from the joint covariance. Pre-event coefficients therefore average exactly zero
by construction, and August 2024 receives its own estimate like any other month. Joint Wald tests on
the pre-event coefficients are invariant to this re-expression.

\begin{keybox}{Why the reference matters, and why it changed}
An earlier version used August 2024 alone as the reference. Against weights fitted on 2022--23, treated
employment and labor force had fallen roughly 0.6--0.9~\lp\ in June--August 2024, \emph{before} the
storm. The same dip appears in the 61 \emph{untreated} North Carolina counties measured against the same
control, so it is a statewide LAUS movement that out-of-state donors did not reproduce, not a storm
precursor; QCEW does not show it. Referencing August meant subtracting that dip from every post-storm
month, understating losses by about 0.6~\lp\ and displacing all 31 pre-event coefficients, which is
exactly what the pre-trend tests then rejected. The pre-event mean is the natural zero once the weights
are fitted with an intercept over the same window.
\end{keybox}

\subsection{Aggregate results}
\begin{center}\small
\begin{tabular}{llrrr}
\toprule
Outcome & Window & Estimate & 95\% CI & $p$\\
\midrule
Resident employment (\lp) & Oct--Dec 2024 & $-1.93$ & $[-3.25,\,-0.62]$ & 0.004\\
 & Jan--Sep 2025 & $-1.19$ & $[-2.51,\,0.13]$ & 0.077\\
 & Nov 2025--Mar 2026 & $-0.10$ & $[-1.87,\,1.67]$ & 0.909\\
Labor force (\lp) & Oct--Dec 2024 & $-1.19$ & $[-2.37,\,-0.01]$ & 0.049\\
 & Jan--Sep 2025 & $-1.02$ & $[-2.31,\,0.26]$ & 0.117\\
 & Nov 2025--Mar 2026 & $-0.59$ & $[-2.24,\,1.06]$ & 0.483\\
Workplace jobs (\lp) & Oct--Dec 2024 & $-1.22$ & $[-2.93,\,0.49]$ & 0.159\\
 & Jan--Sep 2025 & $-0.68$ & $[-3.99,\,2.62]$ & 0.685\\
 & Nov 2025--Mar 2026 & $+0.34$ & $[-3.38,\,4.07]$ & 0.857\\
Unemployment rate (pp) & Oct--Dec 2024 & $+0.70$ & $[0.16,\,1.24]$ & 0.011\\
 & Jan--Sep 2025 & $+0.16$ & $[-0.08,\,0.40]$ & 0.187\\
 & Nov 2025--Mar 2026 & $-0.46$ & $[-0.67,\,-0.26]$ & $<0.001$\\
Jobs / employment (\lp) & Oct--Dec 2024 & $+0.71$ & $[-0.59,\,2.01]$ & 0.280\\
Employment / labor force (\lp) & Oct--Dec 2024 & $-0.75$ & $[-1.32,\,-0.17]$ & 0.011\\
 & Nov 2025--Mar 2026 & $+0.49$ & $[0.27,\,0.70]$ & $<0.001$\\
\bottomrule
\end{tabular}
\end{center}

\section{Damage terciles: the dose design}\label{sec:dose}

\subsection{Construction}
The 39 treated counties are sorted into terciles of exposure, 13 per group, and each tercile receives
its own event-study series against the same synthetic control:
\begin{equation}
y_{it}=\alpha_i+\gamma_t+\sum_{g=1}^{3}\sum_{k\in\mathcal{K}}\beta_{g,k}\,\ind\{i\in\Tr_g\}\ind\{k_t=k\}+\varepsilon_{it},
\end{equation}
with the same sample, weights $\omega_i$, fixed effects, clustering and pre-event-mean recentring as
\eqref{eq:twfe}. The high-minus-low contrast is taken from the joint covariance.

HDSI terciles (cutoffs 49.2 and 78.4) are primary. Rainfall terciles (cutoffs 163\,mm and 256\,mm) are
a robustness comparison. The two classifications agree for 29 of 39 counties (74\%), so rainfall is a
genuine check rather than a restatement.

\subsection{Identification of the dose effects}
Each $\tau_W(g)$ is an average treatment effect for exposure group $g$ relative to never-treated
donors, and therefore requires only the parallel-trends assumption of the main design applied within
that group. The high-minus-low \emph{contrast} is a comparison of differently exposed groups: it is
informative about gradient but is not a randomised dose, because exposure correlates with terrain,
income and industry mix. The paper should state this once and not hedge it repeatedly.

\subsection{Results, HDSI terciles}
\begin{center}\small
\begin{tabular}{llrrrr}
\toprule
Outcome & Window & Low & Mid & High & High $-$ Low\\
\midrule
Resident employment (\lp) & Oct--Dec 2024 & $-0.01$ & $-1.53^{*}$ & $-4.27^{***}$ & $-4.26^{***}$\\
 & Jan--Sep 2025 & $-0.39$ & $-1.40$ & $-1.78$ & $-1.39$\\
 & Nov 2025--Mar 2026 & $+1.16$ & $-0.79$ & $-0.68$ & $-1.84$\\
Labor force (\lp) & Oct--Dec 2024 & $-0.22$ & $-1.59^{*}$ & $-1.76^{*}$ & $-1.53$\\
 & Nov 2025--Mar 2026 & $+0.59$ & $-1.28$ & $-1.08$ & $-1.67$\\
Unemployment rate (pp) & Oct--Dec 2024 & $-0.20^{*}$ & $-0.05$ & $+2.37^{***}$ & $+2.57^{***}$\\
 & Jan--Sep 2025 & $-0.09$ & $-0.07$ & $+0.64^{**}$ & $+0.74^{***}$\\
 & Nov 2025--Mar 2026 & $-0.54^{***}$ & $-0.47^{***}$ & $-0.39^{**}$ & $+0.15$\\
\bottomrule
\end{tabular}
\end{center}
{\footnotesize $^{*}p<0.05$, $^{**}p<0.01$, $^{***}p<0.001$, unadjusted.}

The pattern is monotone in the first quarter and the least-damaged third is a clean null
($-0.01$, $p=0.99$). The aggregate $-1.93$ is therefore not ``the storm cost 2\% of jobs'': it is
thirteen counties losing 4.3\% averaged with thirteen losing nothing.

\subsection{Results, rainfall terciles}
Employment, October--December 2024: low $-0.09$ ($p=0.92$), mid $-2.34$ ($p=0.001$), high $-3.37$
($p=0.001$), contrast $-3.28$ ($p=0.004$). Unemployment rate contrast $+2.13$ ($p<0.001$). The ordering
and the null in the least-exposed third reproduce under a measure containing no FEMA claim behaviour.

\section{Migration as an outcome}

\subsection{Why it cannot be a control}
Migration is realised after treatment. Conditioning on it would be a bad control: it changes the
estimand from the total effect to a non-causal direct effect, and conditioning on a collider opens
paths through unobservables affecting both migration and employment. It is also self-defeating, since
migration is one of the channels under study. The same point applies to HDSI in a different form: being
time-invariant, it is absorbed by the county fixed effects and can only enter interacted with event
time, which is a heterogeneity analysis, not a control.

\subsection{Specification}
Census Vintage 2025 components give annual July-to-July population change and net domestic migration
per 1{,}000 prior residents. With county and year effects, treated-by-year indicators by damage
tercile, 2024 the omitted last unexposed year, and clustering by county:
\begin{equation}
m_{i\,\text{yr}}=\alpha_i+\gamma_{\text{yr}}+\sum_{g=1}^{3}\sum_{\text{yr}\ne2024}\delta_{g,\text{yr}}\,\ind\{i\in\Tr_g\}\ind\{\text{year}=\text{yr}\}+\varepsilon_{i\,\text{yr}} .
\end{equation}
Donors carry the labor synthetic-control weights so that migration is comparable with the labor
outcomes; treated counties carry $1/39$, as in \eqref{eq:twfe}.

\subsection{Choice of comparison scheme}
Both an unweighted comparison (all eligible counties equally) and the synthetic-weighted scheme were
estimated. The pre-tests discriminate decisively. Under the unweighted scheme the low and mid terciles
fail their joint 2021--2023 pre-tests ($p=0.002$ and $0.020$) and the pooled pre-test is borderline
($p=0.056$). Under synthetic weighting all three terciles pass ($p=0.41$, $0.72$, $0.59$) and the
pooled pre-test is clean ($p=0.87$). The synthetic scheme is therefore primary; the unweighted result
is a footnote.

\subsection{Results}
\begin{center}\small
\begin{tabular}{llrrr}
\toprule
Outcome & Group & 2025 effect & 95\% CI & $p$ \\
\midrule
Net domestic migration per 1{,}000 & all 39 & $-3.53$ & $[-7.09,\,0.03]$ & 0.052\\
 & low HDSI & $-1.73$ & $[-5.87,\,2.41]$ & 0.411\\
 & mid HDSI & $-1.54$ & $[-5.28,\,2.19]$ & 0.416\\
 & \textbf{high HDSI} & $\mathbf{-7.32}$ & $[-12.32,\,-2.32]$ & \textbf{0.004}\\
 & high $-$ low & $-5.59$ & $[-10.61,\,-0.58]$ & 0.029\\
Population growth (pp) & high HDSI & $-0.71$ & $[-1.23,\,-0.19]$ & 0.008\\
\bottomrule
\end{tabular}
\end{center}

\subsection{How this closes the mechanism}
In the late window the high-damage tercile's employment has returned toward the counterfactual
($-0.68$, $p=0.56$) while its labor force has not ($-1.08$). Over the year to July 2025 those same
counties run 7.3 net domestic migrants per 1{,}000 \emph{below} the comparison, and 0.71pp below it in
population growth. The orders of magnitude line up: a population-growth shortfall of roughly 0.7pp is
the same size as the labor-force shortfall of roughly 1.1~\lp, so the migration margin can account for
most of the missing labor force in exactly the counties where employment fell hardest. The unemployment
rate ending \emph{below} counterfactual in every tercile is the same fact seen from the denominator.

\begin{plainbox}
\textbf{These counties did not depopulate, and the paper must not say they did.} Net domestic migration
stays \emph{positive} throughout: the most-damaged tercile averages $+11.1$ per 1{,}000 in 2024 and
$+4.8$ in 2025, and its population still grows, by 0.11\%. Only 2 of the 39 treated counties record true
net outflow in 2025 and only 6 record falling population. Every migration statement in this paper is
therefore a statement about a shortfall relative to the counterfactual, not an absolute loss: these are
fast-growing Appalachian counties that kept growing, much more slowly.

This also leaves the direction of the margin unidentified. A net flow cannot separate residents leaving
from prospective arrivals staying away, and because the flow remains positive, the reading that the
storm deterred in-movers is at least as consistent with the data as the reading that it displaced
residents. The two carry different welfare and policy implications and the data here cannot choose
between them.
\end{plainbox}

Two further caveats: migration is annual and covers July 2024--July 2025, only ten post-storm months;
and movers need not be labor-force representative, so the accounting is indicative rather than exact.

\section{Validation}\label{sec:validation}

\subsection{The circularity to avoid}
The weights are chosen to minimise the pre-event gap, so reporting that the pre-event gap is small is
reporting the value of the objective at its own minimum. A check is informative only if it perturbs
something the weights did not see: the information set, the identity of the treated units, or the dose.

\subsection{Check 1: backdated out-of-sample forecast}
Weights refitted on 2022--23 only, predicting January--August 2024:
\begin{center}\small
\begin{tabular}{lcccc}
\toprule
 & Employment & Labor force & Jobs & Unemp.\ rate\\
\midrule
Training RMSE (2022--23) & 0.108 & 0.098 & 0.365 & 0.050\\
Forecast RMSE (Jan--Aug 2024) & 0.501 & 0.594 & 0.593 & 0.101\\
\textbf{Ratio} & \textbf{4.66} & \textbf{6.08} & 1.63 & 2.00\\
\bottomrule
\end{tabular}
\end{center}
A ratio near one would mean donors chosen on two years kept tracking the treated counties in the third.
They did not, for the two LAUS outcomes. \textbf{This is the main weakness of the design and should be
stated before anyone asks.} It is also the reason the definitive weights use the full pre-event window:
with 2024 inside the fit the gap is flat, but that flatness is in-sample, and this table is its price.

\subsection{Check 2: pre-event path}
Two distinct nulls, frequently conflated.

\emph{Does the gap drift?} Linear-trend contrast and equality of half-year bin means: $p\ge0.84$ and
$p\ge0.64$ respectively for all six outcomes. No individual pre-event coefficient is significant for
the four fitted outcomes (four of 32 for the jobs/employment ratio). Within HDSI terciles the drift
tests also pass for employment ($p=0.84$, $0.37$, $0.49$) with zero significant pre-event coefficients
in any group. \emph{Passes.}

\emph{Is the gap exactly zero every month?} Joint Wald test on all 31 pre-event coefficients: rejects
for every outcome, $p$ from $2\times10^{-4}$ to below $10^{-11}$. \emph{Fails} --- but at deviations
with RMSE $0.0014$ in log units for employment against an immediate effect of $0.019$, roughly
one-tenth the size. With 154 clusters and small residual variance this test detects deviations an order
of magnitude below the effect. The correct response is to report both and compare magnitudes, not to
select the favourable test.

\subsection{Check 3: geographic placebo}
Draw 39 counties at random from the 1{,}150 eligible donors, remove them from their own donor pool,
refit the entire procedure (same objective, same $\lambda$, same 32-month window, same intercepts), and
compute the same October--December 2024 shock. 199 draws, all converged.
\begin{center}\small
\begin{tabular}{lccccc}
\toprule
 & Observed & Placebo mean & Placebo SD & $z$ & Most extreme of 199\\
\midrule
Resident employment & $-1.93$ & $+0.07$ & 0.27 & $-7.3$ & $-0.88$\\
Labor force & $-1.19$ & $+0.07$ & 0.26 & $-4.8$ & $-0.78$\\
Workplace jobs & $-1.22$ & $+0.11$ & 0.36 & $-3.7$ & $-0.74$\\
Unemployment rate & $+0.70$ & $+0.00$ & 0.06 & $+11.5$ & $+0.21$\\
\bottomrule
\end{tabular}
\end{center}
The observed shock is not merely the most extreme of 200 values: it is 1.5 to 3.4 times larger than the
most extreme value the procedure produced on untreated counties.

Two presentational traps. First, the saved summary reports \texttt{left\_rank}$\,=1.000$ for the
unemployment rate; this is not a failure but the consequence of the observed shock sitting at the
\emph{top} of the distribution --- read \texttt{right\_rank}$\,=0.005$. Second, the treated aggregate's
scaled pre-fit error is $0.0167$ against a placebo mean of $0.0031$: western North Carolina is about
five times harder to fit than a random group of donor counties, which is why relative fit screens were
non-binding and have been removed.

\subsection{Check 4 (free): the within-treatment placebo}
The least-damaged tercile shows no effect at any horizon. Those thirteen counties are FEMA-positive but
barely hit, they pass through the identical pipeline, and they produce nothing. For many readers this
is more persuasive than the 199 refits, and it costs no additional estimation.

\subsection{What none of these establish}
Parallel trends concerns the \emph{post}-event counterfactual, which is never observed; every check
above concerns the pre-period or untreated units. The geographic placebo is not randomisation
inference, because treatment was not randomly assigned over 39-county groups --- a rank of $1/200$ is a
benchmark of what the procedure yields on untreated data, not a $p$-value for a sharp null. And passing
a low-powered pre-trend test is weak evidence (Roth, 2022), which is why magnitudes accompany every
$p$-value here.

\section{Inference and multiplicity}\label{sec:multiplicity}

Standard errors are clustered by county throughout (154 clusters; $\mathrm{df}=153$). This omits
cross-county regional dependence and treats $\widehat w$ as fixed; the placebo refits partially address
both, since each draw re-estimates its own weights.

Six outcomes $\times$ three windows form an 18-test family. \textbf{Unadjusted $p$-values are reported.}
Holm-adjusted values are computed and retained in the output tables but are not displayed, for two
reasons.

First, the tests are far from independent. Two of the six outcomes are exact linear combinations of the
others --- $\hat\beta^{\text{share}}=\hat\beta^{\text{emp}}-\hat\beta^{\text{lf}}$ holds to machine
precision, since OLS is linear in the dependent variable and the design matrix is common --- so the
correlation matrix of pre-event coefficients has rank 4, not 6, with two exactly zero eigenvalues.
Employment and labor force correlate at 0.90; the unemployment rate and the employment-to-labor-force
ratio at $-0.98$. Holm is valid under arbitrary dependence and therefore prices the worst case, which
is far from this one.

Second, and more importantly, the correction never binds on what the paper claims. The tercile results
carry $p=1.6\times10^{-7}$ (high tercile employment) and $1.1\times10^{-5}$ (high-minus-low contrast);
multiplied by 18 they remain below $10^{-3}$. The compositional late-window results
($p\approx1.5\times10^{-5}$) likewise survive. What does \emph{not} survive adjustment is the aggregate
first-quarter effect ($p=0.0042$, Bonferroni $0.076$) --- and the paper's argument is precisely that the
average is the wrong object. The recommended wording:

\begin{plainbox}
We report unadjusted $p$-values throughout. The two findings the paper rests on --- the damage gradient
and the compositional recovery --- are significant at the 0.1\% level and remain so under Bonferroni
correction across all primary tests. The aggregate first-quarter effect is significant only without
adjustment; that the average is the fragile estimate is itself part of the argument.
\end{plainbox}

A sharper alternative, if a referee presses, is Romano--Wolf stepdown, which bootstraps the joint
distribution of the test statistics and so exploits the actual correlation while still controlling the
family-wise error rate. It would be uniformly more powerful than Holm here.

Finally, displayed intervals are individual county-clustered $t$ intervals, not simultaneous bands.

\section{Exploratory extensions}

\subsection{Longer LAUS horizon}
LAUS outcomes continue to September 2026 (QCEW cannot, ending in March 2026). Over November
2025--September 2026 the unemployment rate sits $0.56$pp below counterfactual ($p<0.001$) and the
employment-to-labor-force ratio $+0.59$~\lp\ ($p<0.001$), while employment ($-0.12$~\lp, $p=0.89$) and
labor force ($-0.72$~\lp, $p=0.40$) are statistically indistinguishable from it. The compositional
reading persists at the longest horizon available. These windows sit outside the primary family.

\section{Design decisions, and what was rejected}

Documenting rejected alternatives is worth the space, because each is a question a reader will ask.

\begin{description}[leftmargin=1em,style=nextline]
\item[Why not plain DiD against all US or all NC counties?] Pre-trends fail. Against uniform donor
weights the treated employment gap moves by $\pm1.5$~\lp\ within 2023 alone; against fitted weights it
moves by 0.14. Plain DiD would attribute ordinary western-NC seasonality and the 2023 tourism rebound
to the storm.

\item[Why not classic synthetic control?] Classic SC matches levels with no intercept and no
regularisation. With 1{,}150 donors and a large diversified treated region, exact level matching is
unnecessary (fixed effects absorb levels) and harmful (it loads weight onto donors sharing a level
rather than a trajectory). The intercept plus ridge is the SDID remedy.

\item[Why terciles rather than a median split?] Tested directly. The median split pools mid-damage
counties into both halves and dilutes from both ends: the employment contrast falls from $-4.26$ to
$-2.87$ and its $p$-value worsens from $1.1\times10^{-5}$ to $5.0\times10^{-4}$ despite larger groups.
It also destroys the clean null --- $-0.54$ ($p=0.45$) instead of $-0.01$ ($p=0.99$). Three points also
exhibit monotonicity; two points show only a difference.

\item[Why not a continuous dose interaction?] A per-SD slope was estimated and has been removed. It is
a causal response only under a stronger parallel-trends assumption than the binned contrasts require
(Callaway, Goodman-Bacon and Sant'Anna, 2024, on continuous treatments), its pre-event joint tests
reject, and nothing downstream depended on it. Binned $\tau_W(g)$ against never-treated donors needs
only the assumption the main design already makes.

\item[Why were the Gaussian exposure curves dropped?] Local-linear kernel curves over exposure were
estimated and removed in November 2026. They described the same gradient the terciles describe, cost a
separate set of assumptions (bandwidth, bootstrap, effective local support of 9--20 counties), and
required roughly 500 words to explain and defend.

\item[Why not control for migration or damage?] See Section~6.1. Both are post-treatment or absorbed;
both belong on the left-hand side.

\item[Why keep six outcomes when two are linear combinations?] They answer different substantive
questions (the commuting margin, the participation margin) and the ratios are the outcomes carrying the
compositional result. Dropping an outcome after observing it is null would be a different and less
defensible act than reporting fewer of them in a short paper.
\end{description}

\section{Limitations}

\begin{enumerate}[leftmargin=1.4em]
\item \textbf{The design is post hoc.} Earlier specifications were inspected before the current one was
frozen, and the November 2026 revisions followed seeing the 2024 diagnostics. This is not
pre-registration. Mitigating: $\lambda$ selection, donor pool, treated roster, outcomes and windows were
not altered; what changed was the fit window and the reference, each documented with its reason.

\item \textbf{The flat pre-period is in-sample.} The definitive weights are fitted over the whole
pre-event window. The backdated test shows donors chosen on 2022--23 forecast 2024 poorly for the LAUS
outcomes. The honest statement is that the control tracks western NC's movements when allowed to see 32
months, and that this cannot be verified out of sample for the post-storm period.

\item \textbf{LAUS is a model, not a survey.} County LAUS estimates are produced from CES, QCEW, UI
claims and CPS inputs via a signal-plus-noise model with statewide benchmarking; the June--August 2024
statewide dip is one visible consequence. Post-storm LAUS may also respond mechanically to UI claims.
QCEW is an administrative census of covered employment and does not share this issue, though it is
noisier at county level.

\item \textbf{Clustering ignores regional dependence and weight uncertainty.} County clusters treat
Wyoming counties as independent of one another and treat $\widehat w$ as fixed.

\item \textbf{Spillovers to donors} are assumed away by geography. A national construction-labor pull
toward North Carolina would bias effects toward zero.

\item \textbf{Treated heterogeneity} is large, which is why the paper reports terciles. Equal weighting
means the estimand is dominated in number by small rural counties.

\item \textbf{The dose contrast is not a randomised dose.} Damage correlates with terrain, income and
industry mix, so high-minus-low compares differently exposed groups.
\end{enumerate}

\section{Anticipated questions}

\begin{description}[leftmargin=1em,style=nextline]
\item[How do you know it was the storm?] Four strands. Timing: a sharp break in the first post-storm
month with no pre-trend. Dose: within the treated group, losses scale with damage, and that gradient is
absent in the pre-event placebo window. Placebos: no random 39-county group of unaffected counties
produced a shock remotely this large under the same procedure. Coherence: employment, labor force, jobs
and unemployment move together in the directions a destruction shock predicts, and workplace and
resident measures fall together rather than diverging.

\item[Your strict pre-tests reject. Is that fatal?] They reject exact month-by-month equality at
deviations of roughly one-tenth the effect size, with 154 clusters. The drift tests do not reject, no
monthly coefficient is individually significant, and the least-damaged tercile is a null. A test that
rejects at any nonzero deviation informs about precision, not about bias of a magnitude that would
overturn the result.

\item[Why is unemployment below counterfactual by 2026?] Two mechanisms consistent with the data and
not separable here: labor-force exit and out-migration (the labor force remains below counterfactual
while employment recovers, mechanically lowering the rate), and reconstruction demand. The
employment-to-labor-force result and the migration estimates are the same fact from different angles.

\item[Can you speak to aggregate regional losses?] Not from this estimand. Buncombe alone is roughly a
third of the treated population; a population-weighted estimate would be close to a study of Asheville.
The equal-county estimand answers a different question, and arguably the more relevant one for policy
toward rural counties.

\item[Why 199 placebo draws?] It makes the smallest attainable rank $1/200=0.005$, which is sufficient
given the observed shock exceeds the most extreme draw by a factor of 1.5--3.4. The constraint is
compute, not inference: each draw refits the full optimisation.
\end{description}

\appendix
\section{Notation}
\begin{center}\small
\begin{tabular}{ll}
\toprule
$\Tr$, $|\Tr|=39$ & treated NC counties (frozen roster); $\Tr_g$ its damage terciles, 13 each\\
$\Dn$, $J=1{,}150$ & eligible donors; 115 carry positive weight after fitting\\
$y^{(k)}_{it}$, $k=1\dots4$ & $100\ln$ employment, $100\ln$ labor force, $100\ln$ jobs, unemployment rate (pp)\\
$\tilde y$ & 2022-centred version\\
$s_k$ & 2022 scale of the treated aggregate; floors 0.25 / 0.15\\
$T_{\text{pre}}$ & Jan 2022--Aug 2024, 32 months (definitive); Jan 2022--Dec 2023 (backdated)\\
$\lambda=1$ & ridge toward uniform, chosen by fit-2022 / predict-2023\\
$\widehat w$, $\omega_i$ & fitted donor weights; regression weights ($1/39$ treated, $\widehat w_i$ donors)\\
$t_0$ & September 2024, excluded transition month\\
$\tilde\beta_k$, $\widehat\tau_W$ & event coefficient and window effect, relative to the pre-event mean\\
$G=154$, $\mathrm{df}=153$ & clusters and degrees of freedom\\
\bottomrule
\end{tabular}
\end{center}

\section{Where each number lives}
\begin{center}\small
\begin{tabular}{>{\raggedright}p{5.6cm}>{\raggedright\arraybackslash}p{9.2cm}}
\toprule
Donor weights, tuning, solver & \texttt{outputs/labor/tables/donor\_weights.csv}, \texttt{inner\_2023\_cv.csv}, \texttt{solver\_audit.csv}, \texttt{training\_2022\_scales.csv}\\
Treated and synthetic paths & \texttt{monthly\_paths.csv}\\
Backdated out-of-sample audit & \texttt{pretrend\_train\_holdout\_ALL.csv} (\texttt{bd\_*} columns)\\
Aggregate event study & \texttt{event\_coefficients.csv}, \texttt{period\_effects\_WITH\_HOLM.csv}\\
Pre-trend tests & \texttt{twfe\_pretrend\_tests.csv}\\
Damage terciles & \texttt{dose\_event\_coefficients.csv}, \texttt{dose\_period\_effects.csv}, \texttt{dose\_pretrend\_tests.csv}, \texttt{dose\_treated\_exposure.csv}\\
Geographic placebos & \texttt{outputs/labor/placebos/}\\
Migration & \texttt{outputs/migration/migration\_by\_damage.csv}\\
Readable report and atlas & \texttt{outputs/reports/}\\
Code & \texttt{src/preprocessing/01--07}, \texttt{src/labor/steps/01--05, 07}, \texttt{src/migration/01--02}, \texttt{src/reporting/}\\
\bottomrule
\end{tabular}
\end{center}

\section*{References}
\begin{itemize}[nosep,leftmargin=1.2em]
\item Abadie, A., Diamond, A., Hainmueller, J. (2010). Synthetic control methods for comparative case studies. \emph{JASA} 105(490).
\item Abadie, A., L'Hour, J. (2021). A penalized synthetic control estimator for disaggregated data. \emph{JASA} 116(536).
\item Angrist, J., Pischke, J.-S. (2009). \emph{Mostly Harmless Econometrics}, \S3.2.3 on bad controls.
\item Arkhangelsky, D., Athey, S., Hirshberg, D., Imbens, G., Wager, S. (2021). Synthetic difference-in-differences. \emph{AER} 111(12).
\item Callaway, B., Goodman-Bacon, A., Sant'Anna, P. (2024). Difference-in-differences with a continuous treatment. NBER WP 32117.
\item Holm, S. (1979). A simple sequentially rejective multiple test procedure. \emph{Scand.\ J.\ Statist.} 6(2).
\item Romano, J., Wolf, M. (2005). Exact and approximate stepdown methods for multiple hypothesis testing. \emph{JASA} 100(469).
\item Roth, J. (2022). Pretest with caution: event-study estimates after testing for parallel trends. \emph{AER: Insights} 4(3).
\end{itemize}

\end{document}
